\begin{tabularx}{\textwidth}{L{31mm}X X}
\toprule
\textbf{Szempont} & \textbf{Túlterhelés} & \textbf{Felüldefiniálás} \\
\midrule
Angol név & Overloading. & Overriding. \\
Kapcsolat & Egy osztályon belül vagy öröklési láncban is előfordulhat. & Öröklődéshez kötődik: a leszármazott újra megadja az ős metódusát. \\
Metódusnév & Azonos. & Azonos. \\
Paraméterlista & Kötelezően eltér. & Azonos szignatúra szükséges. \\
Visszatérési típus & Önmagában nem elég a megkülönböztetéshez. & Lehet azonos vagy kovariáns visszatérési típus. \\
Kiválasztás ideje & Fordítási időben, a referencia statikus típusa és az argumentumok alapján. & Futási időben, az objektum dinamikus típusa alapján. \\
Kulcsszó & Nincs külön kulcsszó. & Java-ban ajánlott az \code{@Override} annotáció használata. \\
Példa & \code{print(String)} és \code{print(int)}. & \code{Dog.makeSound()} felüldefiniálja az \code{Animal.makeSound()} metódust. \\
\bottomrule
\end{tabularx}
